\chapter{Conclusiones y Líneas Futuras}
\label{chap:conclusiones}

\section{Conclusiones generales}
\label{sec:conclusiones_gen}

En este Trabajo Fin de Máster se ha desarrollado una arquitectura integral para el control reactivo y la orquestación de una flota de robots cuadrúpedos orientada a operaciones de búsqueda y rescate en el marco del proyecto RASTREO del Centro de Automática y Robótica (CAR CSIC-UPM).

Entre los principales hitos alcanzados destacan:
\begin{itemize}
    \item La sustitución exitosa de la capa de control original por un driver de bajo nivel en C++20 de alto rendimiento y determinista, reduciendo el consumo de cómputo en más de un 95\%.
    \item La articulación de una arquitectura contenerizada en microservicios desacoplados que aísla de forma segura el hardware crítico frente a posibles fallos de navegación.
    \item La formalización de un árbol de comportamiento tolerante a contingencias y estructurado para la ejecución modular de seis misiones operativas clave.
\end{itemize}

\section{Líneas de trabajo futuro}
\label{sec:lineas_futuras}

Como continuación natural a los desarrollos presentados, se proponen las siguientes líneas de trabajo:
\begin{enumerate}
    \item Despliegue y validación en suelo de las misiones de navegación cooperativa Nav2 con los tres perros físicos de forma simultánea.
    \item Incorporación de algoritmos de percepción visual mediante aprendizaje profundo a bordo de la Raspberry Pi CM5 para el reconocimiento automático de víctimas.
    \item Integración de la brigada canina con plataformas humanoides y estaciones de control de campo dentro del ecosistema completo de RASTREO.
\end{enumerate}
